\documentclass{article}
\usepackage[UTF8]{ctex}
\usepackage{amsmath,amssymb,amsfonts}
\usepackage{enumitem}
\usepackage{geometry}
\geometry{a4paper, margin=1in}
% === 错题本排版优化（不参与题目标识） ===
\setlength{\parindent}{0pt}
\setlength{\parskip}{0.35em plus 0.10em minus 0.08em}
\setlength{\emergencystretch}{3em}
\setlist{itemsep=0.25em, topsep=0.40em, parsep=0pt}
\allowdisplaybreaks[2]
% === 错题本排版优化结束 ===
\begin{document}
\section*{山东大学2026年高等代数真题}
\section*{115．山东大学2026年线性代数与常微分方程试卷}
高等代数部分\\

\subsection{题目1}
1．已知线性方程组

\[
\left\{\begin{array}{l}
\lambda x_{1}+x_{2}+x_{3}=-3 \\
x_{1}+\lambda x_{2}+x_{3}=-3 \\
x_{1}+x_{2}+\lambda x_{3}=-3
\end{array}\right.
\]

当 \(\lambda\) 为何值时，方程组无解，有解，有解时求解。\\

\subsection{题目2}
2．已知 \(A\) 为 \(n\) 阶方阵。\\
（1）若 \(r(A)=1\) ，证明：\(A\) 可以写作 \(\left(\begin{array}{c}a_{1} \\ a_{2} \\ \vdots \\ a_{n}\end{array}\right)\left(\begin{array}{llll}b_{1} & b_{2} & \cdots & b_{n}\end{array}\right)\) ．\\
（2）若 \(r(A)=1\) ，证明：存在常数 \(k\) ，使得 \(A^{2}=k A\) ．\\
（3）若 \(r(A)=1\) ，且存在 \(m \geq 2\) 满足 \(A^{m}=O\) ，证明：\(A^{2}=O\) 。\\

\subsection{题目3}
3．已知 \(A=\left(\begin{array}{cc}A_{11} & A_{12} \\ A_{21} & A_{22}\end{array}\right)\) 为对称矩阵，\(A_{11}\) 为方阵．证明下面几个条件等价：\\
（i）\(A\) 为正定矩阵．\\
（ii）\(A_{11}\) 和 \(A_{22}-A_{21} A_{11}^{-1} A_{12}\) 为正定矩阵．\\
（iii）\(A_{22}\) 和 \(A_{11}-A_{12} A_{22}^{-1} A_{21}\) 为正定矩阵．\\

\subsection{题目4}
4．已知实矩阵

常微分方程部分\\

\subsection{题目5}
2 。求解 \(y=x\left(y^{\prime}+\sqrt{1+\left(y^{\prime}\right)^{2}}\right)\)

\[
A=\left(\begin{array}{cccccc}
a & -b & & & & \\
b & a & 1 & & & \\
& & a & -b & & \\
& & b & a & 1 & \\
& & & a & -b \\
& \text { 初值 } & (1)=1 . & & b & a
\end{array}\right)
\]

其中 \(b \neq 0\) 。求 \(A\) 的不变因子，初等因子以及若尔当标准形。\\

\subsection{题目6}
3．在䣄艮空間 \(\bar{V}\) 中，\(\alpha\) ，其臣零向量，连结番组 \(\alpha_{1}, \alpha_{2}, \cdots, \alpha_{m}\) 。证提：\(y_{\text {两足 }} y(x)\) 在 \((-\infty,+\infty)\) 存在有界解，并求出该解。\\

\subsection{题目7}
4。已知 \(y^{\prime \prime}-y^{\prime}-2 y, \overline{\bar{\alpha}}_{i} f(x)\) 的有界解为,\(\left.h\right),(x)\) ，证明：若 \(\left(x_{i}\right.\), 是周期为 \(T, m\), 的连续函数，则 \(\varphi\) 也是周期为 \(T\) 的连续函数．\\
证眼解 \({ }_{1}, \alpha_{2}, \cdots, \alpha_{m}\) 线性无关。\\

\subsection{题目8}
6．设 \(\alpha_{1}, \alpha_{2}, \cdots, \alpha_{n}\) 为线性空间 \(V\) 的一组基，\(V\) 上的线性变换 \(\sigma\) 满足

\[
\sigma\left(\alpha_{i}\right)=\alpha\left\{\begin{array}{c}
x^{\prime}=-4 x-2 y+\frac{1}{e^{t}-1} \\
+y\left(i=6 y, 2,3 y, n \frac{21}{e^{t}-1}, \sigma\left(\alpha_{n}\right)=0 .\right.
\end{array}\right.
\]

（这）里求 \(x, \sigma_{y}\) 在其 \(\sigma_{1}, ~ \sigma_{\text {的函数 }} n\) 下的矩阵。\\
（2）证明：\(\sigma^{n}=0, \sigma^{n-1} \neq 0\) ，这里 0 为零变换．
\end{document}
